\subsection{The special Cohen series and the limits of square-class Hecke recurrences}

This is a bounded attempt at the same-discriminant unit problem.
The first result strengthens the carried fixed-weight congruence to
the full Fourier series.  The second identifies precisely what the
coefficient recurrences fail to supply.  Neither result asserts a
simultaneous unit or removes a rational class component.

\begin{proposition}[KT-COHEN-UP-REFORMULATION]
Let $p\geq31$ be prime, $p\equiv3\pmod4$, and $b=(p+3)/2$.
Use the Cohen series
$H_r=\sum_{n\geq0}H(r,n)q^n$, and define coefficient extraction by
\[
 U_p\!\left(\sum_{n\geq0}a(n)q^n\right)
                  =\sum_{n\geq0}a(pn)q^n.
\]
Then every coefficient of $H_b$ and $H_2$ is $p$-integral, and
\[
                           H_b\equiv U_p H_2\pmod p
                  \quad\hbox{in }\mathbf F_p[[q]].       \tag{HS.1}
\]
This is a Fourier-series identity.  It does not assert that $U_p$
preserves the level-four, weight-$5/2$ space.
\end{proposition}

\begin{proof}
For a fundamental discriminant $D$, allowing $D=1$ with the
trivial character, put
\[
 P_r(D,m)=\sum_{v\mid m}\mu(v)\chi_D(v)v^{r-1}
                                    \sigma_{2r-1}(m/v).
                                                               \tag{HS.2}
\]
The full Cohen formula is
\[
 H(r,n)=-\beta_r(\chi_D)P_r(D,m),
                  \qquad (-1)^rn=Dm^2.               \tag{HS.3}
\]
At the constant coefficient,
$H(b,0)=-B_{p+3}/(p+3)\equiv-B_4/4=1/120=H(2,0)$.
If $n\equiv1,2\pmod4$, both sides of (HS.1) are zero,
since $pn\equiv3,2\pmod4$ respectively.

Now write $-n=Dm^2$ with $D<0$ fundamental.  The primitive
quadratic character associated with $\chi_D\chi_{-p}$ has positive
fundamental discriminant
\[
 D^\dagger=
 \begin{cases}
  -pD,&p\nmid D,\\
  -D/p,&p\mid D.
 \end{cases}
 \qquad
 m^\dagger=
 \begin{cases}
  m,&p\nmid D,\\
  pm,&p\mid D,
 \end{cases}
 \qquad pn=D^\dagger(m^\dagger)^2.                    \tag{HS.4}
\]
Here the second case includes $D=-p$, $D^\dagger=1$.
The factorization of fundamental discriminants is valid also
when $D$ is even, because $-p$ is odd and its prime occurs in
$D$ to exponent exactly one.

For every integer $a$ prime to the two conductors,
Euler's criterion gives
\[
 \chi_D(a)a^{b-1}
       \equiv\chi_D(a)\chi_{-p}(a)a
       =\chi_{D^\dagger}(a)a\pmod p.                 \tag{HS.5}
\]
Apply the change-of-character Kummer congruence in Guerzhoy's
Proposition~2 to $(\chi_D,b)$ and $(\chi_{D^\dagger},2)$.
Both conductors divide $\Delta p$ for the same $\Delta$ prime to
$p$.  A common regularizer can be chosen with unit residue.
If $\Delta>1$, its quadratic prime-to-$p$ factor is nontrivial;
choose $c\equiv1\pmod p$ and with that factor equal to $-1$.
Then both factors $1-\chi(c)c^r$ are $2$ modulo $p$.
If $\Delta=1$, necessarily $D=-p$, and $c=2$ gives residue
$1-c^2=-3$, again a unit.  In both cases the Euler factors
$1-\chi(p)p^{r-1}$ are $1$ modulo $p$.
The integral measure in Proposition~1 therefore gives
integrality and
\[
                    \beta_b(\chi_D)
                    \equiv\beta_2(\chi_{D^\dagger})\pmod p.
                                                               \tag{HS.6}
\]
This explicitly includes the conductor-divisible-by-$p$ branch,
rather than inferring it from a prime-to-$p$ formula.
The ordinary value $\beta_2(1)=B_2/2=1/12$ is integral as well.

It remains to compare the integer factors (HS.2).
They are multiplicative in $m$.  At a prime $\ell$ their factor is
\[
 P_r(D,\ell^s)=
 \sum_{i=0}^{s}\ell^{i(2r-1)}
   -\chi_D(\ell)\ell^{r-1}
                \sum_{i=0}^{s-1}\ell^{i(2r-1)},       \tag{HS.7}
\]
with an empty second sum when $s=0$.
For $\ell\ne p$ one has
\[
 \ell^{2b-1}\equiv\ell^3,\qquad
 \chi_D(\ell)\ell^{b-1}
      \equiv\chi_{D^\dagger}(\ell)\ell\pmod p.
                                                               \tag{HS.8}
\]
For $\ell=p$, (HS.7) is $1$ modulo $p$ for every $s\geq0$
and every $r\geq2$.  Thus the possible extra factor $p$ in
$m^\dagger$ has no effect, and
$P_b(D,m)\equiv P_2(D^\dagger,m^\dagger)\pmod p$.
Equations (HS.3) and (HS.6) prove (HS.1).

For completeness, the same measure argument gives the
$p$-integrality of every quadratic value $\beta_2$ involved:
its integrand $\chi(a)a$ has a regularizing unit since
$p\geq31$; the exceptional possibility of a residue character
$a^{-2}$ is impossible for a quadratic character at such $p$.
Alternatively, the degree-two Bernoulli-polynomial formula
gives it directly.  Together with (HS.2), this covers all
coefficients of $H_2$; the already checked Cohen integrality
also supplies the coefficients of $H_b$.
\end{proof}

\begin{proposition}[KT-HECKE-SQUARE-CLASS-SCOPE]
Let $r\geq3$ be odd.  Reduction of the full Cohen formula
has the following consequences.

First, if $\beta_r(\chi_D)=0$ modulo $p$, then
$H(r,|D|m^2)=0$ modulo $p$ for every $m$.
In particular, an existing first-unit seed at a discriminant
where $\beta_3$ vanishes cannot be repaired by replacing its
index by square multiples.

Second, the standard Fourier-coefficient equations
\[
 a(\ell^2n)+\left(\frac{-n}{\ell}\right)\ell^{r-1}a(n)
       +\ell^{2r-1}a(n/\ell^2)
       =(1+\ell^{2r-1})a(n),\qquad \ell\nmid2p,         \tag{HS.9}
\]
where $a(n/\ell^2)=0$ if $\ell^2\nmid n$,
do not connect distinct fundamental discriminants.
Even when imposed for every such $\ell$, they leave their
fundamental coefficients formally independent, and they leave
the constant coefficient independent of all positive indices.
\end{proposition}

\begin{proof}
The first assertion is immediate from (HS.2)--(HS.3).
For the second, prescribe arbitrary residues $A_D$ indexed by
negative fundamental discriminants, and an arbitrary constant $c$.
Define a formal series by
\[
 a(0)=c,\qquad
 a(|D|m^2)=A_D P_r(D,m),
 \qquad a(n)=0\ \hbox{at the forbidden plus indices}.   \tag{HS.10}
\]
Every positive allowed index has a unique such factorization.
At $\ell\nmid m$, the $s=1$ case of (HS.7) gives (HS.9),
including $\ell\mid D$ when the character value is zero.
At $\ell\mid m$, the character in (HS.9) is zero and
(HS.7) gives the elementary recurrence
\[
 P_r(D,\ell^{s+1})+\ell^{2r-1}P_r(D,\ell^{s-1})
       =(1+\ell^{2r-1})P_r(D,\ell^s),\qquad s\geq1.
                                                               \tag{HS.11}
\]
The factors at the other primes do not change.
The equation at $n=0$ is simply
$c+\ell^{2r-1}c=(1+\ell^{2r-1})c$.
This proves all the claimed formal equations without any
restriction on the $A_D$.
The Fourier derivative $\Theta=q\,d/dq$ also preserves every
zero coefficient: it multiplies $a(n)$ by $n$.
Iterating it, or combining it with these square-index
coefficient operations, does not create a coefficient in a
previously absent fundamental square class.

For an actual retained example, $p=83$, $D=-43$ satisfies
$\beta_3(\chi_D)=166\equiv0$ and
$\beta_{43}(\chi_D)\equiv76\ne0$.
In (HS.10) with $r=43$, choose $c=1/120$,
$A_{-43}=-76$, and all other $A_D=0$.
This series satisfies every equation (HS.9), has the required
nonzero constant, and has zero coefficientwise product with
$H_3$ at every positive index.  It is a formal example only:
no modularity is asserted.  It isolates the insufficiency of
the Hecke equations themselves, without claiming that the
actual Cohen series has this support.
\end{proof}

\paragraph{Remaining modular question.}
The simultaneous field problem would follow from a nonzero
coefficientwise product
$H(2,pd)H(3,d)$ at an odd fundamental $d<p$.
The first-unit theorem gives a nonzero factor $H(2,pd)$ in that
range, but does not give the other factor.
Neither ordinary multiplication of modular forms nor its
no-zero-divisor property applies to a coefficientwise product.
The modularity of (HS.10) is exactly the omitted global constraint;
the formal construction is not a counterexample to a possible
special modular-form theorem.
The zero set of $\beta_3$ has not been shown to contain only
finitely many fundamental square classes, so a theorem excluding
finite square-class support cannot be invoked here.
No fixed-level modular realization or bounded-weight Sturm
criterion for this coefficientwise product has been established
in this attempt.  Consequently (HS.1) and (HS.9) do not supply
the missing simultaneous unit, even though they give an exact
fixed-weight reformulation and identify the limits of
square-class Hecke amplification.

\paragraph{Inputs and check scope.}
The full Cohen formula is the already checked coefficient
formula, also stated in P.~Guerzhoy,
\emph{Quadratic congruences for Cohen--Eisenstein series},
Glasgow Math.~J.~41 (1999), 141--144, equation~(1), p.~142.
The square-index convention $D=1$ is explicit in
A.~Balog, W.~J.~McGraw and K.~Ono,
\emph{The arithmetic of the coefficients of half-integral weight
Eisenstein series}, Acta Arith.~107 (2003), 337--344,
definitions~(1.2)--(1.4), printed pp.~337--338.
Their formulas at $r=1$, $k=2$ give
$H(2,m^2)=-P_2(1,m)/12$, including $H(2,1)=-1/12$.
Only these definitions are used; their later main theorem
is not needed.
Guerzhoy's Propositions~1--2 and equations~(2)--(4), pp.~143--144,
give precisely the integral measure and character-changing
congruence used in (HS.5)--(HS.6), with conductor divisible
by $p$ permitted.
These statements and the short deduction of Proposition~2 were
read in full.  The published construction of the Mazur measure
remains an imported input.
The recurrence (HS.9) is proved directly from (HS.7) here,
so it requires no unverified converse from a formal eigenseries
to a modular form.
The $p=83$ values are carried exact controls; no prime census,
new discriminant search, or new finite $K$-degree computation
is involved.
